% [Claude] 2026-10-28: placeholder for the motion-coefficient comparison (user request); grey boxes stand in for the images.
% [Claude] 2026-11-05 (user requests): apple (seed 44), Balloon1 (seed 43) and teddy (seed 43) at the first training frame (0_00000), spin (seed 42)
% at training frame 0_00270 (18 s into the 15 fps render). Coefficients = that frame of tests/render_coefs.py's renders (outputs/bdf/modes_vis/
% {som/<S>/bases_10,split_init/<S>/bases_<B>}[_s<seed>]/videos/epoch_0499/motion_coefs.mp4, right half); input = the dataset's training image,
% cropped like the render (Balloon1: 458 x 240 of 470 x 248; DyCheck 720 x 960, uncropped).
\begin{figure}[t]
  \centering
  \newcommand{\coefrow}[3]{\rotatebox{90}{\makebox[#2\linewidth][c]{\footnotesize #3}} & \coefimg{coef_#1_input} & \coefimg{coef_#1_som10} &
    \coefimg{coef_#1_ours} \\}
  \newcommand{\coefimg}[1]{\includegraphics[width=0.3\linewidth]{img/#1}}
  \setlength{\tabcolsep}{1pt}
  \begin{tabular}{@{}c@{\hspace{2pt}}ccc@{}}
    \coefrow{apple_s44_f0}{0.4}{Apple}
    \coefrow{Balloon1_s43_f0}{0.157}{Balloon1}
    \coefrow{teddy_s43_f0}{0.4}{Teddy}
    \coefrow{spin_s42_f270}{0.4}{Spin}
     & \makebox[0.3\linewidth][c]{\footnotesize Input} &
       \makebox[0.3\linewidth][c]{\footnotesize SoM-10} &
       \makebox[0.3\linewidth][c]{\footnotesize Ours} \\
  \end{tabular}
  \caption{\claudetag\ Motion coefficients of Shape of Motion~\cite{wang2025shape} with its default 10 bases and of our method
    ($B$ from its stopping rule: 5 on apple, 4 on Balloon1, 8 on teddy, 9 on spin), rendered from a training view: the first one for
    apple (DyCheck, seed 44), Balloon1 (NVIDIA, seed 43) and teddy (DyCheck, seed 43), frame 270 for spin (DyCheck, seed 42). Each
    foreground Gaussian is coloured by its blending weights: every basis of a model has its own hue, spaced evenly over that model's
    bases, and mixed weights mix the hues; the background is grey. Hues are not matched between the two methods: what matters is which
    parts share one colour.}
  \label{fig:coef-comparison}
\end{figure}
